\section{Local properties of multiplier ideals}
\label{sec:psh-local-ideals}

Multiplier ideals compare singularities through integrability. Their local valuative characterization and strong openness, recalled in \cite{SGPT}, make this comparison compatible with maxima, addition and increasing limits.

\begin{theorem}\label{thm:local-i-inputs}
For germs of plurisubharmonic functions, the following facts hold.
\begin{enumerate}
\item Equality of multiplier ideals at every positive weight is equivalent to equality of the Lelong numbers of the pullbacks at every point on every local modification. Inclusion of the ideals at every positive weight is equivalent to the reverse inequalities for these Lelong numbers.
\item If $p_j$ is an increasing sequence of plurisubharmonic functions converging almost everywhere to a plurisubharmonic function $p$, then, for each point $x$ and each $\lambda>0$,
\[
 \mathcal I(\lambda p_j)_x=\mathcal I(\lambda p)_x
 \quad\text{for all sufficiently large }j.
\]
\end{enumerate}
\end{theorem}

In the first item one works over a coordinate neighborhood of the germ, allowing smaller neighborhoods as necessary. No assertion that every local modification extends to a modification of the compact manifold is needed. The second item is the increasing-limit form of strong openness used in \cite{SGPT}. Smooth additions do not change multiplier ideals or Lelong numbers. Thus these inputs apply to quasi-plurisubharmonic functions after the local correction of \cref{prop:local-correction}.

\begin{lemma}\label{lem:hermitian-i-calculus}
Let $u,u',v,v'$ be quasi-plurisubharmonic functions on $X$. Suppose $u\sim_{\mathcal I}v$ and $u'\sim_{\mathcal I}v'$. Then
\[
 \max(u,u')\sim_{\mathcal I}\max(v,v'),\qquad
 u+u'\sim_{\mathcal I}v+v',\qquad
 a u\sim_{\mathcal I}a v\quad(a>0).
\]
These statements concern quasi-plurisubharmonic functions and require no background to be closed.
\end{lemma}
\begin{proof}
Fix a coordinate neighborhood $U\subset X$ and a smooth local modification $\pi:Y\to U$. Write $\nu(w,y)$ for the Lelong number of a quasi-plurisubharmonic function $w$ at $y$. At every $y\in Y$, the elementary Lelong identities give
\[
 \begin{aligned}
 \nu(\max(u,u')\circ\pi,y)
   &=\min\{\nu(u\circ\pi,y),\nu(u'\circ\pi,y)\},\\
 \nu((u+u')\circ\pi,y)
   &=\nu(u\circ\pi,y)+\nu(u'\circ\pi,y),\\
 \nu(a u\circ\pi,y)&=a\nu(u\circ\pi,y).
 \end{aligned}
\]
The corresponding expressions for $v,v'$ are equal by the first local input. Apply the converse direction of that input on each neighborhood. The Lelong identities are the usual local maximum, addition and scaling identities, unchanged by smooth corrections.
\end{proof}


Let $v,w$ be local quasi-plurisubharmonic functions and let $Y$ be a complex submanifold on which $v|_Y$ is genuine. Restriction and subadditivity give the two inclusions recalled in \cite{SGPT} and \cite[Section 15]{DemAM11}:
\begin{equation}\label{eq:trace-local-ideal-inputs}
 \mathcal I(v|_Y)\subseteq\mathcal I(v)\mathcal O_Y,
 \qquad
 \mathcal I(v+w)\subseteq\mathcal I(v)\mathcal I(w).
\end{equation}
The first follows from local Ohsawa--Takegoshi extension. For the second, weighted holomorphic functions on a product form the Hilbert tensor product of the two weighted spaces; its generators give the product multiplier ideal. Restriction to the diagonal then gives the inclusion. These are local results for psh functions. Smooth corrections extend them to quasi-plurisubharmonic functions with the weight convention of this book. No compact K\"ahler hypothesis is involved.


For a holomorphic submersion, multiplier ideals commute with pullback. Extension of an ideal means the ideal generated by its pullbacks.

\begin{lemma}\label{lem:submersion-multiplier-pullback}
If $X,M$ are complex manifolds, $p:X\to M$ is a holomorphic submersion and $U\in\QPSH(M)$, then for every $t>0$,
\[
 \mathcal I(tU\circ p)
       =\mathcal I(tU)\mathcal O_X.
\]
\end{lemma}
\begin{proof}
Work in product coordinates $(z,\zeta)$ with $p(z,\zeta)=z$. Write $\alpha$ for a multi-index in the fiber variables and $\mathrm{d}V_\zeta$ for their Euclidean volume. A germ in the ideal on the right is a finite sum of bounded holomorphic multiples of pullbacks of germs integrable against $\mathrm{e}^{-tU(z)}$. Fubini proves its weighted integrability on a sufficiently small product neighborhood.

Conversely, let $F(z,\zeta)$ be a weighted-integrable holomorphic germ and expand it in the fiber variables,
\[
 F(z,\zeta)=\sum_\alpha f_\alpha(z)\zeta^\alpha.
\]
The coefficient estimate for a holomorphic function on a polydisc gives, on smaller product neighborhoods,
\[
 |f_\alpha(z)|^2\leq C_\alpha
       \int |F(z,\zeta)|^2\,\mathrm{d}V_\zeta.
\]
Multiply by $\mathrm{e}^{-tU(z)}$ and integrate. Each coefficient germ belongs to $\mathcal I(tU)$. Let $J$ be the extension of this ideal to the local analytic ring $R$ of the product. Every finite Taylor sum belongs to $J$, and the remainder of total fiber degree at least $N$ belongs to $(\zeta)^N$. Thus the class of $F$ in $R/J$ lies in $(\zeta)^N(R/J)$ for every $N$. The Krull intersection theorem for the finite module $R/J$ over the Noetherian local analytic ring gives intersection zero. Hence $F\in J$. This is the usual local algebra form of convergence for coherent ideals. Smooth corrections and smooth positive volume densities do not change the integrability ideals.
\end{proof}
